\documentclass[oneside,a4paper,14pt]{extarticle}

\usepackage[russian]{babel}
\usepackage{fontspec}
\setmainfont{Times New Roman}
\setmonofont{Courier New}
\usepackage[a4paper,top=2cm,bottom=2cm,left=2.5cm,right=1.5cm]{geometry}
\usepackage{textcomp}
\usepackage{indentfirst}
\usepackage{listings}
\usepackage{xcolor}
\usepackage{enumitem}
\usepackage{caption}
\usepackage{amsmath}
\usepackage{float}
\usepackage{titlesec}
\usepackage{hyperref}

\titleformat{\section}{\normalsize\bfseries}{\thesection}{1em}{}
\titleformat{\subsection}{\normalsize\bfseries}{\thesubsection}{1em}{}

\renewcommand\baselinestretch{1.33}\normalsize
\setlength{\parindent}{1.25cm}

\renewcommand{\lstlistingname}{\,{}}
\lstset{
	basicstyle=\ttfamily\footnotesize,
	keywordstyle=\color{blue},
	commentstyle=\color{green!60!black},
	stringstyle=\color{red},
	numbers=left,
	numberstyle=\tiny\color{gray},
	frame=single,
	breaklines=true,
	showstringspaces=false,
	tabsize=2,
	xleftmargin=1em
}

\begin{document}

\newpage
\thispagestyle{empty}
\begin{center}
	МИНИСТЕРСТВО НАУКИ И ВЫСШЕГО ОБРАЗОВАНИЯ\\
	РОССИЙСКОЙ ФЕДЕРАЦИИ
	ФЕДЕРАЛЬНОЕ ГОСУДАРСТВЕННОЕ БЮДЖЕТНОЕ\\
	ОБРАЗОВАТЕЛЬНОЕ
	УЧРЕЖДЕНИЕ ВЫСШЕГО ОБРАЗОВАНИЯ\\
	«ВЯТСКИЙ ГОСУДАРСТВЕННЫЙ УНИВЕРСИТЕТ»\\
	Институт математики и информационных систем\\
	Факультет автоматики и вычислительной техники\\
	Кафедра электронных вычислительных машин
\end{center}

\vspace{20mm}

\begin{center}
	Отчёт по лабораторной работе №6\\
	по дисциплине\\
	<<Программирование>>
\end{center}

\vspace{48mm}

\noindent
Выполнил студент гр. ИВТб-1303-06-00 \hfill \rule[-0,5mm]{30mm}{0.15mm}\,/Шубин В. Ю./
\newline
Проверил преподаватель кафедры ЭВМ \hfill \rule[-0,5mm]{30mm}{0.15mm}\,/Баташев П.А./

\vfill

\begin{center}
	Киров\\
	2026
\end{center}

\newpage
\thispagestyle{empty}

\section*{Цель работы}
Освоить методы векторизованной предобработки данных с использованием библиотеки NumPy. Получить навыки работы со структурированными массивами, булевыми масками, групповой агрегацией и созданием производных признаков для задач промышленного мониторинга.

\section*{Постановка задачи}
Необходимо выполнить предобработку данных GPS-телеметрии транспорта (вариант~8). Исходный файл \texttt{data\_variant\_8.csv} содержит 2~000~000 записей и следующие поля:
\begin{itemize}[leftmargin=*]
	\item \texttt{ts (int32)} --- время телеметрии;
	\item \texttt{vehicle\_id (int16)} --- идентификатор автомобиля;
	\item \texttt{speed (float32)} --- скорость, км/ч;
	\item \texttt{fuel (float32)} --- уровень топлива, \%;
	\item \texttt{engine (float32)} --- обороты двигателя, тыс/мин;
	\item \texttt{acc (uint8)} --- статус ускорения (0/1).
\end{itemize}

Правила варианта: удалить строки со \texttt{speed < 0} или \texttt{fuel > 100}; заменить \texttt{fuel < 0} на~0; ограничить \texttt{engine} 85-м перцентилем. Группировка выполняется по \texttt{vehicle\_id} со расчётом среднего \texttt{fuel} и максимального \texttt{speed}. Для окна $k=45$ строится скользящее среднее \texttt{fuel\_consumption}, а разность соседних значений \texttt{speed} даёт оценку ускорения. Отдельный акцент сделан на обнаружении резкого скачка уровня топлива --- сброса датчика или заправки.

Требуется реализовать 10 этапов обработки:
\begin{enumerate}[leftmargin=*]
	\item загрузка данных в структурированный массив \texttt{numpy.dtype}, оценка объёма памяти, подсчёт NaN/Inf;
	\item векторизованная фильтрация аномалий и очистка с обрезкой выбросов;
	\item группировка по \texttt{vehicle\_id}, расчёт статистик и Z-score нормализация \texttt{fuel} внутри групп;
	\item скользящее среднее для \texttt{fuel\_consumption} с окном $k=45$ через кумулятивную сумму, вычисление разности \texttt{speed};
	\item создание производных признаков: \texttt{speed\_per\_engine} и \texttt{refuel\_jump};
	\item условная агрегация по группам для записей с резким скачком топлива;
	\item лаговые признаки (lag=1) и анализ распределения знаков разницы;
	\item робастная замена выбросов по методу IQR внутри каждой группы;
	\item проверка согласованности данных по предметным правилам и замена некорректных значений;
	\item частотный анализ поля \texttt{vehicle\_id}, объединение редких категорий и оформление отчёта.
\end{enumerate}

Ограничение: запрещено использование циклов \texttt{for} по строкам исходного массива. Все операции выполняются векторизованно через булевы маски, \texttt{np.where}, \texttt{np.clip}, \texttt{np.unique}; цикл допустим только по уникальным группам.

\section*{Описание реализации}
Обработка выполнена в среде Jupyter с использованием библиотеки NumPy. Код организован в виде последовательных ячеек ноутбука \texttt{lab6.ipynb}, каждая из которых соответствует одному этапу лабораторной работы. Промежуточные массивы сохраняются в \texttt{*.npy}, итоговые метрики --- в \texttt{metrics\_lab6.json}.

\section*{Этап 1: Загрузка и подготовка типов}
Для чтения файла использована функция \texttt{np.genfromtxt} с заданным структурированным типом данных:
\begin{lstlisting}[language=Python]
dtype = np.dtype([
    ("ts", np.int32), ("vehicle_id", np.int16),
    ("speed", np.float32), ("fuel", np.float32),
    ("engine", np.float32), ("acc", np.uint8)
])
data = np.genfromtxt("data_variant_8.csv", delimiter=",",
                     skip_header=1, dtype=dtype, encoding="utf-8")
\end{lstlisting}
Объём памяти рассчитан через атрибут \texttt{nbytes}: исходный массив занимает 38\,000\,000 байт (36,24~МиБ). Подсчёт некорректных значений выполнен через \texttt{np.isfinite}: доля NaN/Inf составила 3,00\,\% в среднем (по полям \texttt{speed} и \texttt{fuel} --- 3,01\,\%, что превышает порог 3\,\%).

\section*{Этап 2: Векторизованная очистка}
Нефинитные значения заменены медианой соответствующего поля. Аномалии обнаружены составной булевой маской:
\begin{lstlisting}[language=Python]
anomaly_mask = (clean["speed"] < 0) | (clean["fuel"] > 100)
clean = clean[~anomaly_mask]
clean["fuel"] = np.where(clean["fuel"] < 0, 0, clean["fuel"])
p85 = np.percentile(clean["engine"], 85)
clean["engine"] = np.clip(clean["engine"], None, p85)
\end{lstlisting}
Обнаружено 129\,599 строк с аномалиями (6,48\,\%): \texttt{speed < 0} --- 97\,641 запись, \texttt{fuel > 100} --- 33\,630 записей. После удаления осталось 1\,870\,401 запись. 85-й перцентиль оборотов составил 11,170 тыс/мин.

\section*{Этап 3: Группировка и нормализация}
Группировка реализована через \texttt{np.unique} с циклом по уникальным значениям \texttt{vehicle\_id} (цикл по группам разрешён). Получено 100 групп мощностью от 18\,411 до 19\,029 записей. Для каждой группы рассчитаны среднее, стандартное отклонение \texttt{fuel} и максимальный \texttt{speed}. Z-score нормализация выполнена по формуле
\begin{equation}
	z = \frac{x - \mu_{group}}{\sigma_{group} + \varepsilon}, \qquad \varepsilon = 10^{-8}.
	\label{eq:zscore}
\end{equation}
Поскольку в исходных данных нет явного поля расхода, определено производное значение \texttt{fuel\_consumption = 100 - fuel} --- доля израсходованного топлива в баке.

\section*{Этап 4: Скользящее окно и разность}
Скользящее среднее вычислено за $O(N)$ через кумулятивную сумму; первые $k-1$ значений усредняются по доступной части окна, поэтому искусственных пропусков не возникает:
\begin{lstlisting}[language=Python]
cs = np.cumsum(x, dtype=np.float64)
out[:k-1] = cs[:k-1] / np.arange(1, k)
out[k-1:] = (cs[k-1:] - np.concatenate(([0.0], cs[:-k]))) / k
accel = np.diff(speed, prepend=speed[0])
\end{lstlisting}
Для $k=45$ построен признак \texttt{fuel\_ma45}. Диапазон оценок ускорения \texttt{accel} составил $[-182{,}51;\ 182{,}24]$ км/ч между соседними замерами. Обнаружено 48\,603 резких скачка \texttt{fuel} более чем на 20 п.~п. (2,60\,\%).

\section*{Этапы 5--10}
\begin{itemize}[leftmargin=*]
	\item \textbf{Производные признаки.} \texttt{speed\_per\_engine = speed / (engine + $10^{-3}$)} с защищённым знаменателем и заменой нефинитных значений на~0; \texttt{refuel\_jump = max(diff(fuel), 0)} --- положительный скачок уровня топлива. Положительные значения \texttt{refuel\_jump} встретились в 707\,215 записях.
	\item \textbf{Условная агрегация.} Маска \texttt{refuel\_jump > 20} отбирает 48\,603 записи (2,60\,\%). Для каждого из 100 автомобилей сформирован массив формы $(100,\ 3)$ со столбцами \texttt{[group\_id, mean, median]}; отдельно сохранён 90-й перцентиль. Средний расход на таких участках колеблется от 44,77 до 51,99\,\%, p90 --- от 78,40 до 79,24\,\%.
	\item \textbf{Лаговые признаки.} \texttt{fuel\_lag1 = np.roll(fuel\_consumption, 1)} с сохранением первого значения; разница текущего и лагового значений распределена через \texttt{np.sign} и \texttt{np.unique}.
	\item \textbf{Робастная замена IQR.} Внутри каждой группы вычислены границы $[Q1 - 1{,}5 \cdot IQR,\ Q3 + 1{,}5 \cdot IQR]$; выбросы заменены медианой группы. Заменено 111\,048 значений (5,94\,\%).
	\item \textbf{Проверка логики.} Проверены ограничения $0 \le \texttt{speed} \le 200$, $0 \le \texttt{fuel} \le 100$, $0{,}8 \le \texttt{engine} \le 8$, $\texttt{acc} \in \{0,1\}$. Нарушения обнаружены в 1\,848\,436 записях (98,83\,\%), главным образом из-за того, что генератор выдаёт \texttt{acc} в диапазоне 0--49. Некорректные флаги заменены эталонными значениями через \texttt{np.where}.
	\item \textbf{Частотный анализ.} Из 100 категорий \texttt{vehicle\_id} 49 встречаются менее чем в 1\,\% записей и объединены в \texttt{OTHER} с кодом~0; частые категории получили коды 1--51.
\end{itemize}

\section*{Результаты выполнения}
Обработка выполнена на наборе из 2\,000\,000 записей. Основные метрики представлены в таблице~\ref{tab:metrics}.

\begin{table}[H]
	\centering
	\caption{Результаты предобработки данных GPS-телеметрии}
	\label{tab:metrics}
	\renewcommand{\arraystretch}{1.2}
	\begin{tabular}{|p{4.6cm}|p{2.9cm}|p{8.3cm}|}
		\hline
		Показатель & Значение & Комментарий \\
		\hline
		Объём исходных данных & 36,24 МиБ & 2 млн строк $\times$ 6 полей структуры \\
		\hline
		Итоговый массив & 94,54 МиБ & Добавлено 9 производных полей \\
		\hline
		Доля NaN/Inf & 3,00\,\% & По \texttt{speed} и \texttt{fuel} --- 3,01\,\%, порог превышен \\
		\hline
		Строк с аномалиями & 6,48\,\% & Удалено 129\,599 записей \\
		\hline
		Количество групп & 100 & Мощность 18\,411--19\,029 \\
		\hline
		Записей со скачком топлива & 2,60\,\% & 48\,603 записи, условие \texttt{refuel\_jump > 20} \\
		\hline
		Заменено выбросов (IQR) & 5,94\,\% & 111\,048 значений \\
		\hline
		Нарушений логики & 98,83\,\% & В основном некорректный диапазон \texttt{acc} \\
		\hline
		Редких категорий & 49 из 100 & Объединены в \texttt{OTHER} (код 0) \\
		\hline
	\end{tabular}
\end{table}

Анализ временных сдвигов (lag=1 для \texttt{fuel\_consumption}):
\begin{itemize}[leftmargin=*]
	\item рост значения: 37,79\,\% записей;
	\item падение значения: 37,81\,\% записей;
	\item без изменений: 24,40\,\% записей.
\end{itemize}
Симметричная динамика роста и падения указывает на естественные колебания скорости и расхода при движении транспорта без выраженного тренда.

\section*{Выводы}
В ходе выполнения лабораторной работы были освоены методы векторизованной предобработки данных с использованием библиотеки NumPy. Ключевые результаты:
\begin{enumerate}[leftmargin=*]
	\item реализована полная обработка массива из 2 млн записей без циклов по строкам, что обеспечило высокую производительность;
	\item применены предметно-ориентированные правила очистки: физические ограничения скорости и уровня топлива, обрезка оборотов двигателя по 85-му перцентилю;
	\item групповая нормализация и IQR-замена выбросов позволили учесть индивидуальные особенности автомобилей;
	\item временные признаки (скользящее среднее, лаги, разности) выявили резкие скачки датчика топлива и оценили стиль вождения;
	\item сжатие редких категорий предотвратило разреженность признакового пространства для последующего машинного обучения.
\end{enumerate}

Особенности предобработки данных GPS-телеметрии:
\begin{itemize}[leftmargin=*]
	\item телеметрия содержит значительную долю аппаратного шума: около 3\,\% пропусков и столько же отрицательных значений в каждом вещественном поле;
	\item сброс датчика топлива проявляется как резкий положительный скачок, который уверенно обнаруживается через \texttt{diff(fuel)} и сглаживается окном $k=45$;
	\item группировка по \texttt{vehicle\_id} критически важна: автомобили различаются по расходу и максимальной скорости;
	\item векторизованные операции NumPy позволяют обрабатывать миллионы записей без поэлементных циклов.
\end{itemize}

Итоговый массив \texttt{data\_final.npy} физически непротиворечив, нормализован и готов к решению прикладных задач: прогнозированию расхода топлива, контролю стиля вождения и предиктивному обслуживанию автопарка.

\section*{Ссылка на рабочую среду}
Код лабораторной работы размещён в локальном ноутбуке \texttt{lab6.ipynb} в папке \texttt{progalab 6}; все ячейки выполнены, результаты сохранены в файлах \texttt{data\_*.npy} и \texttt{metrics\_lab6.json}.

\end{document}
